% Figura 4 de la postulación: espacio de decisión y estados operacionales.
\documentclass[border=4pt]{standalone}
% Preámbulo común de las figuras 1 a 4 de la postulación (se carga con \input en cada figura).
% Los cortes y umbrales vienen de umbrales_calibrados.tex (lo escribe figuras_postulacion.py).
\usepackage{fontspec}
\setmainfont{Cambria}
\usepackage[table]{xcolor}
\usepackage{array,booktabs}
\usepackage{tikz}
\usetikzlibrary{arrows.meta,positioning,calc,fit,backgrounds,decorations.pathreplacing}
\definecolor{fondo}{HTML}{F2F1EC}
\definecolor{txt2}{HTML}{52514E}
% estados, ordenados por magnitud: E3 despeje < E4 tormenta < E5 peligro
\definecolor{cE0}{HTML}{D9D8D2}\definecolor{cE1}{HTML}{1BAF7A}\definecolor{cE2}{HTML}{EDA100}
\definecolor{cE3}{HTML}{2A78D6}\definecolor{cE4}{HTML}{E34948}\definecolor{cE5}{HTML}{4A3AA7}
\definecolor{cSV}{HTML}{2A78D6}\definecolor{cH}{HTML}{EB6834}\definecolor{cN}{HTML}{52514E}
% Generado por figuras_postulacion.py desde salidas/sistema/parametros.json. No editar a mano.
\newcommand{\cortesS}{\textbf{0{,}5} & 5 & \textbf{12{,}5} & 20 & \textbf{25}}
\newcommand{\cortesF}{0{,}05 & 0{,}25 & 0{,}50 & 0{,}75 & 0{,}95}
\newcommand{\cortesr}{0{,}5 & 1 & 2 & 3 & 4}
\newcommand{\cortesdur}{1 & 6 & 12 & 18 & 23}
\newcommand{\cortesG}{10 & 15 & 20 & 25 & 30}
\newcommand{\cortesR}{\textbf{0{,}5} & 2{,}5 & 5 & \textbf{10} & 15}
\newcommand{\cortesHNS}{10 & 20 & 30 & 50 & 80}
\newcommand{\cortesRS}{10 & 20 & 40 & 60 & 100}
\newcommand{\pesoS}{0{,}45}
\newcommand{\pesoF}{0{,}15}
\newcommand{\pesor}{0{,}10}
\newcommand{\pesodur}{0{,}10}
\newcommand{\pesoG}{0{,}10}
\newcommand{\pesoR}{0{,}10}
\newcommand{\SVrest}{1}
\newcommand{\SVparc}{2{,}5}
\newcommand{\SVcierre}{3{,}5}
\newcommand{\Hrest}{3}
\newcommand{\Hparc}{4}
\newcommand{\Hcierre}{5}
\newcommand{\Nrest}{5}
\newcommand{\Nparc}{30}
\newcommand{\Ncierre}{50}
\newcommand{\cdesp}{6}
\newcommand{\expAlto}{1}
\newcommand{\expMedio}{0{,}5}
\newcommand{\hnSobrecarga}{30}

\newcolumntype{L}[1]{>{\raggedright\arraybackslash}p{#1}}
\newcolumntype{C}{>{\centering\arraybackslash}p{0.62cm}}
% encabezado de notas con degradado del color del índice
\newcommand{\notas}[1]{\cellcolor{#1!12}1 & \cellcolor{#1!28}2 & \cellcolor{#1!45}3 & \cellcolor{#1!65}\textcolor{white}{4} & \cellcolor{#1!85}\textcolor{white}{5}}
\tikzset{
  tit/.style={anchor=north west, font=\small\bfseries, inner sep=0pt},
  caja/.style={draw=#1, line width=0.9pt, rounded corners=3pt, align=center, minimum width=4.2cm, minimum height=1.05cm, fill=white},
  dato/.style={draw=txt2, dashed, rounded corners=3pt, align=center, minimum width=4.5cm, minimum height=0.85cm},
  est/.style={rounded corners=2pt, align=center, minimum width=1.98cm, minimum height=1.05cm, font=\scriptsize, text=white, inner sep=1pt},
  niv/.style={rounded corners=2pt, align=center, minimum width=2.45cm, minimum height=0.62cm, font=\scriptsize, inner sep=1pt},
  llave/.style={decorate, decoration={brace, amplitude=3pt}, draw=txt2},
  flecha/.style={-{Stealth[length=2mm]}, thick, draw=txt2},
  eti/.style={font=\scriptsize, color=txt2, align=center}}

\begin{document}
\begin{tikzpicture}[font=\footnotesize]
% a) Espacio de decisión
\node[tit] at (0,0) {a) Espacio de decisión};
\begin{scope}[shift={(0.8,-4.75)}, x=0.86cm, y=0.62cm]
  % regiones: el estado es el E más alto entre severidad y peligro (filas H = 0..5)
  \fill[cE0] (0,0) rectangle (5,6);
  \fill[cE1] (1,0) rectangle (5,6);  \fill[cE1] (0,3) rectangle (5,6);
  \fill[cE2] (2.5,0) rectangle (5,6); \fill[cE2] (0,4) rectangle (5,6);
  \fill[cE4] (3.5,0) rectangle (5,6);
  \fill[cE5] (0,5) rectangle (5,6);
  \draw[white, line width=0.4pt] (0,0) grid[xstep=10, ystep=1] (5,6);
  \draw[txt2] (0,0) rectangle (5,6);
  \foreach \x/\l in {0/0, 1/\SVrest, 2.5/\SVparc, 3.5/\SVcierre, 5/5} \draw[txt2] (\x,0) -- (\x,-0.12) node[below, font=\scriptsize, inner sep=1pt] {\l};
  \foreach \y in {0,...,5} \node[left, font=\scriptsize, inner sep=2pt] at (0,\y+0.5) {\y};
  \node[font=\scriptsize, color=txt2] at (2.5,-0.85) {severidad $SV$};
  \node[font=\scriptsize, color=txt2, rotate=90] at (-0.55,3) {peligro $H$};
  \node[font=\scriptsize] at (0.5,1.5) {$E_0$};
  \node[font=\scriptsize, text=white, align=center] at (1.75,1.5) {$E_1$\\restricción};
  \node[font=\scriptsize] at (3.0,1.5) {$E_2$};
  \node[font=\scriptsize, text=white] at (1.25,3.5) {$E_1$ restricción};
  \node[font=\scriptsize] at (1.25,4.5) {$E_2$ cierre parcial};
  \node[font=\scriptsize, text=white] at (2.5,5.5) {$E_5$ cerrado por peligro};
  \node[font=\scriptsize, text=white, rotate=90, align=center] at (4.25,2.5) {$E_4$ cerrado por tormenta};
  % tercer eje: nieve por despejar
  \begin{scope}[shift={(5.75,0)}]
    \fill[cE0] (0,0) rectangle (0.55,0.8);
    \fill[cE1] (0,0.8) rectangle (0.55,2.6);
    \fill[cE2] (0,2.6) rectangle (0.55,4.2);
    \fill[cE3] (0,4.2) rectangle (0.55,6);
    \draw[txt2] (0,0) rectangle (0.55,6);
    \foreach \y/\l in {0/0, 0.8/\Nrest, 2.6/\Nparc, 4.2/\Ncierre} \draw[txt2] (0.55,\y) -- (0.67,\y) node[right, font=\scriptsize, inner sep=1pt] {\l};
    \node[font=\scriptsize, text=white, rotate=90] at (0.275,5.1) {$E_3$};
    \node[font=\scriptsize, color=txt2, rotate=90] at (1.45,3) {nieve por despejar $N$ (cm)};
  \end{scope}
\end{scope}
\node[anchor=north west, font=\scriptsize, text width=7.2cm, align=left] at (0,-5.9)
  {\textbf{Regla:} cada índice tiene tres niveles, y el estado es el $E$ más alto que alcanza cualquiera de ellos. Se reabre cuando el peligro baja de \Hcierre\ y la nieve por despejar baja de \Ncierre\ cm.};

% b) Estados operacionales
\node[tit] at (7.9,0) {b) Estados operacionales, de menor a mayor magnitud};
\node[anchor=north west, inner sep=0pt] at (7.9,-0.45) {\scriptsize\setlength{\tabcolsep}{3pt}\renewcommand{\arraystretch}{1.3}%
\begin{tabular}{>{\raggedright\arraybackslash}p{1.95cm}L{3.55cm}L{2.45cm}}
\toprule
Estado & Qué implica & Se activa con \\
\midrule
\cellcolor{cE0}$E_0$ abierto & operación normal & $SV < \SVrest$, $H < \Hrest$ y $N \le \Nrest$ cm \\
\cellcolor{cE1}\textcolor{white}{$E_1$ abierto con restricción} & tránsito en ambos sentidos, con cadenas obligatorias (alerta preventiva del protocolo) & $SV \ge \SVrest$, $H \ge \Hrest$ o $N > \Nrest$ cm \\
\cellcolor{cE2}$E_2$ cierre parcial & solo bajada; cierre por tramos; suben emergencias y abastecimiento (alerta operativa) & $SV \ge \SVparc$, $H \ge \Hparc$ o $N > \Nparc$ cm \\
\cellcolor{cE3}\textcolor{white}{$E_3$ cerrado por despeje} & ya no nieva; trabajan las máquinas & $N > \Ncierre$ cm tras un cierre \\
\cellcolor{cE4}\textcolor{white}{$E_4$ cerrado por tormenta} & cerrado mientras dura la tormenta (alerta crítica) & $SV \ge \SVcierre$ \\
\cellcolor{cE5}\textcolor{white}{$E_5$ cerrado por peligro} & avalancha o remoción en masa; las cuadrillas no entran & $H = \Hcierre$ \\
\bottomrule
\end{tabular}};
\end{tikzpicture}
\end{document}
